\documentclass{article}
\usepackage{graphicx} % Required for inserting images
\usepackage{biblatex}
\usepackage{xcolor}
\usepackage{physics}
\addbibresource{refs.bib}
\usepackage{geometry}
\usepackage{subcaption}


\RequirePackage{xcolor}
\newcommand{\revision}[1]{{\color{red}{#1}}}
\newcommand{\ts}[1]{{\color{red}- \textbf{Thiago}: #1 - }}
\newcommand{\ca}[1]{{\color{orange} - \textbf{Cesar}: #1 - }}

\geometry{
    left=15mm,
    right=15mm,
    top=15mm,
    bottom=15mm
}
\usepackage{amsmath}
\usepackage{amssymb}

\title{List of simulations}
\author{Cesar Amaral, Thiago T. Tsutsui}
\date{September 2025}

\begin{document}

\maketitle

\section{Only dephasing}

Considering $\kappa=0$ and $\gamma=10^{-3}$ and $g_0=1$.

\begin{itemize}
    \item Wigner Negativity. Time from $t=0$ to $t=15$.
    \begin{itemize}
        \item Gaussian with varying width, $T=7.5$, $\sigma=-1$, $1.8<\zeta<3$. We employed $250$ values of the parameter $\zeta$.
        When plotting the values in a Jupyter notebook employ the 125+1th parameter for $\zeta_{\text{av}}$.
        \item Gaussian with varying peak time, $\zeta=2.4$, $\sigma=-1$, $4.5<T<10.5$. We employed $100$ values of the parameter $T$.
        When plotting the values in a Jupyter notebook employ the 50+1th parameter for $T_{\text{av}}$.
        \item Cosine with varying frequency, $\phi=0$, $0<\omega<2\pi/3$. We employed $100$ values of the parameter $\omega$.
        When plotting the values in a Jupyter notebook employ the 25+1th parameter for $\omega_{\text{min}}$ because the first parameter is $0$, and the 50+1th parameter for $\omega_{\text{av}}$.
    \end{itemize}
    \item Coherence. Time from $t=0$ to $t=50$.
    \begin{itemize}
        \item Gaussian with varying width, $T=25$, $\sigma=-1$, $6<\zeta<10$. We employed $100$ values of the parameter $\zeta$.
        When plotting the values in a Jupyter notebook employ the 50+1th parameter for $\zeta_{\text{av}}$.
        \item Gaussian with varying peak time, $\zeta=8$, $\sigma=-1$, $15<T<35$. We employed $100$ values of the parameter $T$.
        When plotting the values in a Jupyter notebook employ the 50+1th parameter for $T_{\text{av}}$.
        \item Cosine with varying frequency, $\phi=0$, $0<\omega<\pi/5$. We employed $200$ values of the parameter $\omega$.
        When plotting the values in a Jupyter notebook employ the 50+1th parameter for $\omega_{\text{min}}$ because the first parameter is $0$, and the 100+1th parameter for $\omega_{\text{av}}$.
    \end{itemize}
    \item Magic. Time from $t=0$ to $t=50$.
    \begin{itemize}
         \item Gaussian with varying width, $T=25$, $\sigma=-1$, $6<\zeta<10$. We employed $100$ values of the parameter $\zeta$.
        When plotting the values in a Jupyter notebook employ the 50+1th parameter for $\zeta_{\text{av}}$.
        \item Gaussian with varying peak time, $\zeta=8$, $\sigma=-1$, $15<T<35$. We employed $100$ values of the parameter $T$.
        When plotting the values in a Jupyter notebook employ the 50+1th parameter for $T_{\text{av}}$.
        \item Cosine with varying frequency, $\phi=0$, $0<\omega<\pi/5$. We employed $200$ values of the parameter $\omega$.
        When plotting the values in a Jupyter notebook employ the 50+1th parameter for $\omega_{\text{min}}$ because the first parameter is $0$, and the 100+1th parameter for $\omega_{\text{av}}$.
    \end{itemize}
    \item Entanglement. Time from $t=0$ to $t=50$.
    \begin{itemize}
        \item Gaussian with varying width, $T=25$, $\sigma=-1$, $6<\zeta<10$. We employed $100$ values of the parameter $\zeta$.
        When plotting the values in a Jupyter notebook employ the 50+1th parameter for $\zeta_{\text{av}}$.
        \item Gaussian with varying peak time, $\zeta=8$, $\sigma=-1$, $15<T<35$. We employed $100$ values of the parameter $T$.
        When plotting the values in a Jupyter notebook employ the 50+1th parameter for $T_{\text{av}}$.
        \item Cosine with varying frequency, $\phi=0$, $0<\omega<\pi/5$. We employed $200$ values of the parameter $\omega$.
        When plotting the values in a Jupyter notebook employ the 50+1th parameter for $\omega_{\text{min}}$ because the first parameter is $0$, and the 100+1th parameter for $\omega_{\text{av}}$.
    \end{itemize}
\end{itemize}

\section{Only cavity damping}

Considering $\gamma=0$ and $\kappa=10^{-2}$ and $g_0=1$.

\begin{itemize}
    \item Wigner Negativity. Time from $t=0$ to $t=15$.
    \begin{itemize}
        \item Gaussian with varying width, $T=7.5$, $\sigma=-1$, $1.8<\zeta<3$. We employed $250$ values of the parameter $\zeta$.
        When plotting the values in a Jupyter notebook employ the 125+1th parameter for $\zeta_{\text{av}}$.
        \item Gaussian with varying peak time, $\zeta=2.4$, $\sigma=-1$, $4.5<T<10.5$. We employed $100$ values of the parameter $T$.
        When plotting the values in a Jupyter notebook employ the 50+1th parameter for $T_{\text{av}}$.
        \item Cosine with varying frequency, $\phi=0$, $0<\omega<2\pi/3$. We employed $100$ values of the parameter $\omega$.
        When plotting the values in a Jupyter notebook employ the 25+1th parameter for $\omega_{\text{min}}$ because the first parameter is $0$, and the 50+1th parameter for $\omega_{\text{av}}$.
    \end{itemize}
    \item Coherence. Time from $t=0$ to $t=50$.
    \begin{itemize}
        \item Gaussian with varying width, $T=25$, $\sigma=-1$, $6<\zeta<10$. We employed $100$ values of the parameter $\zeta$.
        When plotting the values in a Jupyter notebook employ the 50+1th parameter for $\zeta_{\text{av}}$.
        \item Gaussian with varying peak time, $\zeta=8$, $\sigma=-1$, $15<T<35$. We employed $100$ values of the parameter $T$.
        When plotting the values in a Jupyter notebook employ the 50+1th parameter for $T_{\text{av}}$.
        \item Cosine with varying frequency, $\phi=0$, $0<\omega<\pi/5$. We employed $200$ values of the parameter $\omega$.
        When plotting the values in a Jupyter notebook employ the 50+1th parameter for $\omega_{\text{min}}$ because the first parameter is $0$, and the 100+1th parameter for $\omega_{\text{av}}$.
    \end{itemize}
    \item Magic. Time from $t=0$ to $t=50$.
    \begin{itemize}
         \item Gaussian with varying width, $T=25$, $\sigma=-1$, $6<\zeta<10$. We employed $100$ values of the parameter $\zeta$.
        When plotting the values in a Jupyter notebook employ the 50+1th parameter for $\zeta_{\text{av}}$.
        \item Gaussian with varying peak time, $\zeta=8$, $\sigma=-1$, $15<T<35$. We employed $100$ values of the parameter $T$.
        When plotting the values in a Jupyter notebook employ the 50+1th parameter for $T_{\text{av}}$.
        \item Cosine with varying frequency, $\phi=0$, $0<\omega<\pi/5$. We employed $200$ values of the parameter $\omega$.
        When plotting the values in a Jupyter notebook employ the 50+1th parameter for $\omega_{\text{min}}$ because the first parameter is $0$, and the 100+1th parameter for $\omega_{\text{av}}$.
    \end{itemize}
    \item Entanglement. Time from $t=0$ to $t=50$.
    \begin{itemize}
        \item Gaussian with varying width, $T=25$, $\sigma=-1$, $6<\zeta<10$. We employed $100$ values of the parameter $\zeta$.
        When plotting the values in a Jupyter notebook employ the 50+1th parameter for $\zeta_{\text{av}}$.
        \item Gaussian with varying peak time, $\zeta=8$, $\sigma=-1$, $15<T<35$. We employed $100$ values of the parameter $T$.
        When plotting the values in a Jupyter notebook employ the 50+1th parameter for $T_{\text{av}}$.
        \item Cosine with varying frequency, $\phi=0$, $0<\omega<\pi/5$. We employed $200$ values of the parameter $\omega$.
        When plotting the values in a Jupyter notebook employ the 50+1th parameter for $\omega_{\text{min}}$ because the first parameter is $0$, and the 100+1th parameter for $\omega_{\text{av}}$.
    \end{itemize}
\end{itemize}

\section{Specific parameters}

Considering $\kappa=10^{-2}$ and $\gamma=10^{-3}$ and $g_0=1$.

\begin{itemize}
    \item Wigner Negativity. Time from $t=0$ to $t=15$.
    \begin{itemize}
        \item Gaussian with varying width, $T=7.5$, $\sigma=-1$, $1.8<\zeta<3$. We employed $250$ values of the parameter $\zeta$.
        Simulate specifically for $\zeta_{\text{av}}=2.4$.
        \item Gaussian with varying peak time, $\zeta=2.4$, $\sigma=-1$, $4.5<T<10.5$. We employed $100$ values of the parameter $T$.
       Simulate specifically for $T_{\text{av}}=7.5$.
        \item Cosine with varying frequency, $\phi=0$, $0<\omega<2\pi/3$. We employed $100$ values of the parameter $\omega$.
         Simulate specifically for $\omega_{\text{min}}=\pi/6$ and $\omega_{\text{av}}=\pi/3$.
    \end{itemize}
    \item Coherence. Time from $t=0$ to $t=50$.
    \begin{itemize}
        \item Gaussian with varying width, $T=25$, $\sigma=-1$, $6<\zeta<10$. We employed $100$ values of the parameter $\zeta$.
       Simulate specifically for $\zeta_{\text{av}}=8$.
        \item Gaussian with varying peak time, $\zeta=8$, $\sigma=-1$, $15<T<35$. We employed $100$ values of the parameter $T$.
       Simulate specifically for $T_{\text{av}}=25$.
        \item Cosine with varying frequency, $\phi=0$, $0<\omega<\pi/5$. We employed $200$ values of the parameter $\omega$.
             Simulate specifically for $\omega_{\text{min}}=\pi/20$ and $\omega_{\text{av}}=\pi/10$.
    \end{itemize}
    \item Magic. Time from $t=0$ to $t=50$.
    \begin{itemize}
         \item Gaussian with varying width, $T=25$, $\sigma=-1$, $6<\zeta<10$. We employed $100$ values of the parameter $\zeta$.
       Simulate specifically for $\zeta_{\text{av}}=8$.
        \item Gaussian with varying peak time, $\zeta=8$, $\sigma=-1$, $15<T<35$. We employed $100$ values of the parameter $T$.
         Simulate specifically for $T_{\text{av}}=25$.
        \item Cosine with varying frequency, $\phi=0$, $0<\omega<\pi/5$. We employed $200$ values of the parameter $\omega$.
          Simulate specifically for $\omega_{\text{min}}=\pi/20$ and $\omega_{\text{av}}=\pi/10$.
    \end{itemize}
    \item Entanglement. Time from $t=0$ to $t=50$.
    \begin{itemize}
        \item Gaussian with varying width, $T=25$, $\sigma=-1$, $6<\zeta<10$. We employed $100$ values of the parameter $\zeta$.
         Simulate specifically for $\zeta_{\text{av}}=8$.
        \item Gaussian with varying peak time, $\zeta=8$, $\sigma=-1$, $15<T<35$. We employed $100$ values of the parameter $T$.
        Simulate specifically for $T_{\text{av}}=25$.
        \item Cosine with varying frequency, $\phi=0$, $0<\omega<\pi/5$. We employed $200$ values of the parameter $\omega$.
         Simulate specifically for $\omega_{\text{min}}=\pi/20$ and $\omega_{\text{av}}=\pi/10$.
    \end{itemize}
\end{itemize}



\printbibliography

\end{document}
